\begin{exercise}[The Gibbons--Hawking--York term]{Tong, Cambridge Part III Sheet 3}
Vary the Einstein--Hilbert action \(I_G\) on a region with non-null boundary.
Show which normal derivative of \(\delta g_{\mu\nu}\) remains, and prove that
adding the Gibbons--Hawking--York term gives a well-posed Dirichlet variational
principle.
\end{exercise}

\begin{exercise}[ADM decomposition of the action]{MIT 8.962, Problem Set 10}
Decompose the scalar curvature using lapse, shift, spatial metric, and
extrinsic curvature.  After including the boundary term, derive the ADM form
of \(I_G\).
\end{exercise}

\begin{exercise}[Lapse and shift as Lagrange multipliers]{Cambridge Part III, General Relativity examples}
Vary the ADM action with respect to lapse and shift.  Recover the Hamiltonian
and momentum constraints and explain why lapse and shift carry no conjugate
dynamical momenta.
\end{exercise}

\begin{exercise}[A nonminimally coupled scalar]{Cambridge Part III, 2024 exam}
For
\[
I=\int\sqrt{-g}\left[-\frac12(\nabla\phi)^2
-\frac12\xi R\phi^2-V(\phi)\right]\dd^4x,
\]
derive the scalar equation and metric stress tensor.  Find the conformal value
of \(\xi\) for a massless scalar in four dimensions.
\end{exercise}

\begin{exercise}[Scalar--tensor theory in the Einstein frame]{Tong, Cambridge Part III Sheet 3}
Starting from
\[
I=\int\sqrt{-g}\left[\frac12F(\phi)R
-\frac12Z(\phi)(\nabla\phi)^2-U(\phi)\right]\dd^4x,
\]
perform a conformal transformation that makes the curvature coefficient
constant.  Determine the transformed scalar kinetic coefficient and
potential.
\end{exercise}

\begin{exercise}[Tetrad postulate and spin connection]{McGreevy, Problem Set 6}
Impose \(\nabla_\mu e^a{}_\nu=0\) and solve for the torsion-free spin
connection in terms of the tetrad and Christoffel connection.  Show how it
transforms under a local Lorentz rotation.
\end{exercise}

\begin{exercise}[Gauss--Bonnet as a boundary variation]{Cambridge Part III, General Relativity examples}
Vary the four-dimensional Gauss--Bonnet action and show that its bulk
Euler--Lagrange tensor vanishes identically.  Explain why the same density
becomes dynamical in dimensions greater than four.
\end{exercise}

\begin{exercise}[Noether current for diffeomorphisms]{Cambridge Part III, 2022 exam}
For a generally covariant Lagrangian, vary the fields by a Lie derivative
along \(\xi^\mu\).  Derive the on-shell conserved Noether current and explain
why it can be written locally as the divergence of an antisymmetric charge
two-form.
\end{exercise}

\begin{exercise}[Canonical momentum and boundary stress]{MIT 8.962, Problem Set 9}
After adding the gravitational boundary term, vary the on-shell action with
respect to the induced boundary metric.  Derive the canonical momentum
proportional to \(K^{ij}-Kh^{ij}\) and identify the associated quasilocal
stress tensor.
\end{exercise}
